\paragraph{The metaplectic comparison at two and the extracted character.}
We specify how the preceding index calculation fixes the multiplier,
including its dyadic part. Let
\[
 Q_0=X_1^2+X_2^2+X_3^2,\qquad
 \phi_{0,p}=1_{\mathbb Z_p^3},\qquad
 \theta(\tau)=\sum_{n\in\mathbb Z}e^{2\pi i n^2\tau}.
\]
Use the same additive character and metaplectic lifts for the
rank-three Weil representations of $Q_B$ and $Q_0$. The latter
theta series is literally $\theta(\tau)^3$; this fixes the meaning
of the ordinary theta multiplier cubed, including its eighth roots
of unity. For either representation the nontrivial central element
of the metaplectic cover acts by $-1$. Consequently the ratio of
two eigenvalues computed using the same lift is independent of
that lift. In particular a central factor in multiplying the
three lifts in the displayed triangular factorization cancels
in this ratio. We compare eigenvalues, and do not choose a group
section of the cover over $\Gamma_0(2^{A_2})$.

Here is the comparison on the three factors. Both upper unipotents
act trivially on the indicated vectors. Choose $A_p$ large enough
that both lower unipotents do also. This is the Fourier-support
argument above for $f_p$; for $\phi_{0,p}$ it requires no positive
depth at an odd prime and depth two at two, since the bilinear dual
of $\mathbb Z_2^3$ for $Q_0$ is $\tfrac12\mathbb Z_2^3$ and
$4Q_0(\tfrac12\mathbb Z_2^3)\subset\mathbb Z_2$.
The Weyl scalar cancels in conjugating an upper unipotent to a
lower one, so it introduces no additional eigenvalue. On a
diagonal unit $u$, the scaling part acts by $\alpha_p(u)$ on
$f_p$ and by one on $\phi_{0,p}$. The quotient of the Weil-index
parts, by the preceding calculation, is
\[
 \frac{\gamma(uQ_B)/\gamma(Q_B)}
      {(\gamma(u)/\gamma(1))^3}
       =(u,\det Q_B)_p.                                      \tag{O.5a}
\]
If the inverse Weil-index convention is used for the diagonal
operator, both numerator and denominator are inverted; the
quotient is unchanged because the Hilbert symbol has order two.
Thus this computation does not leave an undetermined inverse
or a dyadic root of unity. These operator comparisons are precisely
the scaled pullback, Fourier conjugation and scalar relations of
\cite[I.13--I.14 and II.25--II.28]{weil}.

Here $\det Q$ means the determinant of the matrix in
$Q(x)=x^{\mathsf T}Ax$. The associated bilinear Gram matrix is
$2A$. If that convention is used instead, both $Q$ and $Q_0$
acquire the same factor $2^3$ in their determinants, which cancels
in the comparison. In particular no extra squareclass of two
is introduced at the dyadic place.

For a general ternary form $Q$ in this argument, at each prescribed
prime take $A_p\ge1$ and increase $A_p$
also to the conductor of $u\mapsto(u,\det Q)_p$. The eigenvalue
relative to the ordinary theta multiplier, on a matrix whose
upper-left entry is $a$, is then
\[
                    \alpha_p(a)(a,\det Q)_p.                  \tag{O.5b}
\]
Indeed the triangular factorization has just the two trivial
unipotent factors and this diagonal factor, and its lift ambiguity
cancels as explained above. This also verifies multiplication:
the upper-left entry of a product is congruent to the product
of the upper-left entries modulo $p^{A_p}$, so both local
characters in (O.5b) multiply. The argument therefore applies
to every matrix in the indicated congruence subgroup at a prescribed
prime, not only to a choice of generators. At an unramified prime
outside the support, the full spherical invariance established
below treats matrices with nonunit upper-left entry as well.

For $Q_B$ the determinant squareclass is one at every place,
including two. Thus (O.5b) has no determinant factor. For the
extracted form $Q_B'=Q_B/m$ it is $m^{-3}\equiv m$.
At an odd prime $p\mid m$ and a unit $u$, and at two, respectively,
the additional factors are explicitly
\[
 (u,m)_p=(u/p)^{v_p(m)},\qquad
 (u,m)_2=(-1)^{((u-1)/2)((m-1)/2)}.
\]
Here $m$ is odd and $m\equiv3\pmod4$, so the second factor is
$(-1)^{(u-1)/2}$. It is already trivial on $1+4\mathbb Z_2$.
At every odd prime outside $m$ the factor on units is one. There
is consequently no new level prime, and there is no omitted
dyadic character.

Finally take a positive integer $d$ prime to the complete level.
The finite operator used in the classical transformation is the
one at $\gamma^{-1}$, with upper-left entry $d$. The root-character
factor is therefore $\prod_p\alpha_p(d)=\eta_B(d)$.
Hilbert reciprocity gives
\[
 \prod_{p\mid 2m}(d,m)_p
   =\prod_{q\mid d}(m/q)^{v_q(d)}=(m/d).                       \tag{O.5c}
\]
Indeed the infinite symbol is one because $m,d>0$, all other
finite symbols are one, and at an odd $q\mid d$ the symbol is
$(m/q)^{v_q(d)}$. This proves the displayed product without
discarding its dyadic factor. Both sides extend to the same
even Dirichlet character. Since $h_*=-m$, division by the genus
character now gives
\[
 \frac{\eta_B(d)(m/d)}{\eta_B(d)\chi_{h_*}(d)}=(-1/d).
\]
Thus the characters in (O.5), with the stated inverse-idele
convention for $\eta_B$, are exactly the characters supplied
by these local Weil operators.

Take $A_p$ also at least the conductor exponent of $\alpha_p$.
For $\gamma=\left(\begin{smallmatrix}a&b\\c&d\end{smallmatrix}\right)$
with $c\in p^{A_p}\mathbb Z_p$ and $a$ a unit, use
\[
 \gamma=n^-(c/a)m(a)n(b/a).
\]
The scalar is the diagonal theta factor times $\alpha_p(a)$.
In passing to the classical transformation law one applies the
finite operator at $\gamma^{-1}$, whose upper-left entry is $d$.
Thus the classical Dirichlet character is
$\eta_B(d)=\prod_p\alpha_p(d)$, not its inverse.  Its value at
$-1$ is one, as required by the weight-$3/2$ theta multiplier.

